% Guide d'utilisation — Thème Proxmox Cyberpunk V29
% Compilation, depuis ce dossier : lualatex guide_theme_proxmox.tex (ou : tectonic guide_theme_proxmox.tex)
% Sous Debian : apt install texlive-luatex texlive-latex-extra texlive-pictures texlive-lang-french
%               fonts-lato fonts-dejavu-core
% Fournis ici : la capture 3_interface.png et la police Orbitron (SIL OFL 1.1, texte dans ../OFL-Orbitron.txt).
\documentclass[9pt]{extarticle}
\usepackage[a4paper,margin=11mm,top=9mm,bottom=8mm]{geometry}
\usepackage{fontspec}
\usepackage[french]{babel}
\usepackage{xcolor,graphicx,tikz,multicol,enumitem,tabularx,array}
\usepackage[most]{tcolorbox}
\usepackage{microtype}
\usepackage{amssymb}
\usepackage[hidelinks]{hyperref}

\setmainfont{Lato}[UprightFont=*-Regular,BoldFont=*-Bold,ItalicFont=*-Italic,BoldItalicFont=*-BoldItalic]
\setmonofont{DejaVu Sans Mono}[Scale=0.84]
\newfontfamily\orbitron{Orbitron-900.ttf}[Path=./]
\newfontfamily\orbitronb{Orbitron-700.ttf}[Path=./]
\newfontfamily\latoblack{Lato-Black}

\definecolor{nuit}{HTML}{0A0E1A}
\definecolor{nuit2}{HTML}{14233F}
\definecolor{cyan}{HTML}{00B8D9}
\definecolor{cyanvif}{HTML}{00E5FF}
\definecolor{vert}{HTML}{00A86B}
\definecolor{rouge}{HTML}{CC2244}
\definecolor{encre}{HTML}{1B2433}
\definecolor{gris}{HTML}{5B6676}
\definecolor{fondcode}{HTML}{0F1A2E}

\graphicspath{{./}{../captures/}}
\pagestyle{empty}
\setlength{\parindent}{0pt}
\setlength{\parskip}{2pt}
\setlength{\columnsep}{7mm}
\color{encre}
\setlist{nosep,leftmargin=11pt,itemsep=1pt,label={\color{cyan}\scriptsize$\blacktriangleright$}}

% Titre de rubrique
\newcommand{\rubrique}[1]{%
  \par\vspace{5pt}{\orbitronb\fontsize{9.6}{11}\selectfont\color{nuit2}\MakeUppercase{#1}}\par
  \vspace{-2pt}{\color{cyan}\rule{\linewidth}{0.9pt}}\par\vspace{2pt}}

% Bloc de commandes (fond sombre comme le thème)
\newtcolorbox{commande}{enhanced,colback=fondcode,colframe=cyan!70!black,boxrule=0.5pt,arc=2pt,
  left=4pt,right=4pt,top=2.5pt,bottom=2.5pt,fontupper=\ttfamily\color{cyanvif!35!white},before skip=3pt,after skip=4pt}
\newcommand{\cmt}[1]{{\color{gris!80!white}\##1}}
\newcommand{\opt}[1]{{\ttfamily\bfseries\color{nuit2}#1}}
\newcommand{\optt}[1]{{\ttfamily\bfseries\color{nuit2}\fontsize{7.6}{9}\selectfont #1}}
\newcommand{\chemin}[1]{{\ttfamily\small #1}}

\begin{document}

% ───────────── Bandeau ─────────────
\begin{tikzpicture}[remember picture,overlay]
  \fill[nuit] (current page.north west) rectangle ([yshift=-35mm]current page.north east);
  \shade[left color=nuit, right color=nuit2] ([yshift=-35mm]current page.north west) rectangle ([yshift=-32mm]current page.north east);
  \fill[cyanvif] ([yshift=-35mm]current page.north west) rectangle ([yshift=-35.8mm]current page.north east);
  \foreach \x/\y/\r in {0.12/0.35/0.5, 0.31/0.7/0.35, 0.55/0.25/0.45, 0.72/0.62/0.3, 0.88/0.4/0.5, 0.95/0.75/0.3}
    \fill[cyanvif,opacity=0.55] ([xshift=\x*\paperwidth,yshift=-\y*26mm]current page.north west) circle (\r pt);
\end{tikzpicture}
\vspace*{-1mm}
{\orbitron\fontsize{22}{24}\selectfont\color{cyanvif}THÈME PROXMOX CYBERPUNK}\par
\vspace{1mm}
{\fontsize{11}{13}\selectfont\color{white}\textbf{Guide d'installation pour collègues} \;—\; V29 \;—\; Proxmox VE 8 et 9}\par
\vspace{0.8mm}
{\small\color{cyanvif!60!white}Damien SCONTRINO — BTS SIO, lycée Louis Armand (Nogent-sur-Marne) — réutilisation pédagogique libre (CC BY-NC-SA 4.0)}\par
\vspace{1mm}
{\small\color{white}Télécharger le script et ce guide :
\href{https://github.com/SCONTRINO-Hub/proxmox-theme-neon-cyberpunk}{\color{cyanvif}\textbf{github.com/SCONTRINO-Hub/proxmox-theme-neon-cyberpunk}}}\par
\vspace{4mm}

\begin{multicols}{2}

\rubrique{Ce que fait ce script}
Il habille l'interface web de Proxmox VE en \textbf{bleu néon} :
\begin{itemize}
  \item fond sombre animé, menus, fenêtres et graphiques assortis ;
  \item icônes colorées selon l'état : {\color{rouge}\textbf{VM arrêtée}}, VM en marche,
        \textcolor{violet}{\textbf{modèle}} ;
  \item un \textbf{badge} dans la barre du haut (établissement, nom du nœud…) ;
  \item en option, désactivées par défaut : message d'accueil après la connexion,
        retrait de la fenêtre « No valid subscription ».
\end{itemize}
\textbf{Aucun téléchargement, aucun service redémarré, retour arrière en une commande.}
C'est un thème \textbf{sombre} : il impose le fond sombre même si le nœud était en interface claire ;
la désinstallation rend la main à ce réglage.

\rubrique{Prérequis}
\begin{itemize}
  \item Proxmox VE \textbf{8.x} ou \textbf{9.x}, n'importe quel nom de nœud.
  \item Un accès \textbf{root} (console, shell web ou SSH). Rien à installer : bash, perl et coreutils suffisent.
  \item En cluster : lancer le script \textbf{sur chaque nœud}.
\end{itemize}

\rubrique{Après une mise à jour de Proxmox}
Quand Proxmox met à jour son interface, le thème disparaît, sans rien casser : \textbf{relancez le script}
(menu \textbf{1} puis Entrée), ou laissez-le \textbf{revenir tout seul} (\textbf{o} à la question du menu,
ou \opt{-{}-auto-maj}).

\columnbreak

\rubrique{Installation}
Copier le script sur le serveur, puis en root :
\begin{commande}
chmod +x install\_theme\_proxmox.sh\\
./install\_theme\_proxmox.sh \cmt{ menu interactif}
\end{commande}
Le menu propose : \textbf{1} installer ou mettre à jour, \textbf{2} désinstaller. Il pose quatre questions
(badge, message, fenêtre d'abonnement, mises à jour). Ensuite : \textbf{Ctrl+F5} dans le navigateur.
Installation directe, sans question (réglages non précisés inchangés) :
\begin{commande}
./install\_theme\_proxmox.sh --installer --badge "\%NODE\% - Lycée X"
\end{commande}



\rubrique{Ce qui est modifié (et rien d'autre)}
\begin{itemize}
  \item \chemin{index.html.tpl} : 2 lignes \texttt{<link>} (la base sombre de Proxmox, puis le thème)
        et, si option, 1 \texttt{<script>} — toutes repérées par la marque \texttt{theme-proxmox-cyberpunk} ;
  \item une feuille de style (police intégrée) et, si option, un petit script ;
  \item \chemin{proxmoxlib.js} \emph{uniquement} avec \opt{--sans-abonnement} ;
  \item avec \opt{--auto-maj} : un fichier dans \chemin{/etc/apt/apt.conf.d/}, une copie du script
        dans \chemin{/usr/local/sbin/}.
\end{itemize}
{\raggedright Originaux copiés dans \chemin{/var/backups/theme-proxmox-cyberpunk/}.
Venant de la V28 : son crochet APT (jamais lancé par apt, par erreur), \chemin{/opt} et son journal
sont remplacés ou retirés ; ses sauvegardes (\chemin{/root/sauvegardes-theme-proxmox/}) restent.\par}

\end{multicols}

\vspace{1mm}
\newcommand{\cadre}[1]{%
  \begin{tcolorbox}[enhanced,colback=nuit,colframe=cyan!70!black,boxrule=0.6pt,arc=2pt,
    left=0pt,right=0pt,top=0pt,bottom=0pt,drop fuzzy shadow=cyan!40!black]
    \includegraphics[width=\linewidth]{#1}
  \end{tcolorbox}}

\begin{center}\begin{minipage}{0.89\linewidth}\cadre{3_interface.png}\end{minipage}\end{center}\vspace{-2pt}
{\small\textbf{Vue d'un nœud.} VM arrêtée en rouge, modèles en violet, conteneur, réseau et
stockages ; badge, menus, grilles et barre des tâches au même style.
\textcolor{gray}{Interface réelle ; seuls les noms de comptes d'étudiants ont été retirés.}}\par


\vspace{1mm}
\begin{tcolorbox}[enhanced,colback=cyan!6,colframe=cyan!60!black,boxrule=0.8pt,arc=3pt,
                  left=8pt,right=8pt,top=4pt,bottom=4pt]
{\orbitron\fontsize{10}{12}\selectfont\color{cyan!55!black}ET SI ÇA NE VOUS PLAÎT PAS ?}\par\vspace{2mm}
\small Tout se défait, et rien ne vous engage : menu \textbf{« 2) Désinstaller »}, ou
\texttt{./install\_theme\_proxmox.sh -{}-desinstaller}. Le script retire ses lignes et ses fichiers,
et remet \chemin{proxmoxlib.js} tel qu'il l'avait trouvé (si un autre outil l'a retouché depuis,
il n'y touche pas et le dit).
\end{tcolorbox}

\vspace{1mm}
{\scriptsize\color{gris}\raggedright CC BY-NC-SA 4.0 — © 2025-2026 Damien SCONTRINO (damien.scontrino@ac-creteil.fr).
Police Orbitron © 2018 The Orbitron Project Authors, SIL OFL 1.1 (\texttt{OFL-Orbitron.txt}).\par}

\end{document}
